\documentclass[11pt]{article}

% --- arXiv format requirements (rubric item B1): ---
%   single-spaced, 10--14pt type, >= 1in margins, machine-readable, no line
%   numbers, no watermarks, no highlighted text, no margin notes.
\usepackage[margin=1in]{geometry}
\usepackage{amsmath,amssymb}
\usepackage{booktabs}
\usepackage{microtype}
\usepackage[hidelinks]{hyperref}
\usepackage{url}

\title{Fast Weights Predict, Not Evaluate: A Hebbian Associative World Model for Interception}

\author{Gabriel Hill\\\texttt{gabriel@familyfungroup.com}}

\date{}

\begin{document}

\maketitle

\begin{abstract}
\noindent
The Dragon Hatchling (BDH) architecture is built on a fast-weight associative memory --- a
matrix of stimulus--response associations maintained by outer-product plasticity and read out
by similarity-weighted recall. We ask what this primitive is \emph{for}: is it a value
function, an evaluative map from states to expected return, or a world model, a predictive map
from states to next states? We argue for the latter and settle the question empirically in a
controlled interception task, where a turn-rate-limited pursuer must catch a moving target.
Three findings emerge. \textbf{(1)} Used as a value function, the memory collapses: a linear
TD learner with the same associative readout reaches $\sim\!2$--$10\%$ of the lead-pursuit
ceiling, and backpropagated DQN and PPO do not exceed the \emph{no-prediction reflex}.
\textbf{(2)} Used as a world model with greedy lead-pursuit planning, the same memory reaches
$\sim\!94\%$ of the lead-pursuit ceiling. \textbf{(3)} On persistently curving prey --- where
intercepting a trajectory, not a point velocity, is what matters --- the learned world model
beats first- and second-order analytic extrapolation by $17\%$ and $193\%$ respectively,
because rolling forward the learned dynamics is more accurate than a truncated Taylor series.
We interpret the results as supporting a Dreamer-style division of labor: fast-weight memories
are best understood as \emph{predictive} substrates trained by dense self-supervised error,
with sparse, outcome-driven planning on top, rather than as value-function approximators.
\end{abstract}

\section{Introduction}\label{sec:intro}

A central question in biologically plausible machine learning is what a particular synaptic
plasticity rule \emph{computes}. The Dragon Hatchling (BDH) architecture
\cite{bdh2025} places at its core a fast-weight associative memory: a weight matrix updated by
an outer product between a pre-synaptic activity vector and a post-synaptic signal, and
queried by an inner product that recalls stored associations weighted by similarity. This is
one of the oldest and most general ideas in connectionism --- the linear associator, the
Hopfield memory, the ``fast-weight programmer'' \cite{schmidhuber1992,schlag2021}. But
\emph{what should the post-synaptic signal encode?} Two readings are possible, and they lead
to very different systems:

\begin{itemize}
\item \textbf{The evaluative reading.} The memory stores state $\to$ \emph{value} associations,
  forming a value function $V(s)$ or $Q(s,a)$ that is learned by bootstrapped
  temporal-difference (TD) updates and queried by a policy to select actions.
\item \textbf{The predictive reading.} The memory stores state $\to$ \emph{next-state}
  associations, forming a world model $\hat{s}_{t+1} = f(s_t)$ that is trained by
  self-supervised prediction error and rolled forward (imagined) to plan.
\end{itemize}

These readings are not cosmetic: they make opposite claims about the role of the same
plasticity rule. The BDH paper's own framing --- and the broader literature on reward-modulated
Hebbian learning \cite{fremaux2016} --- has tended toward the evaluative reading. We argue the
predictive reading is the correct one, and we support the argument with a controlled
experiment.

We choose \textbf{interception} (pursuit) as the test bed because it makes the two readings
sharply distinguishable. Intercepting a moving target requires \emph{predicting where the
target will be}, not merely evaluating how good the current state is. It is also a domain with
clean analytic baselines (first- and second-order extrapolation) against which a learned model
can be judged, and with an oracle ceiling we can compute. Finally, interception is a canonical
biologically relevant behavior, present from dragonflies to predators to human ball-players.

Our contributions are:
\begin{enumerate}
\item A \textbf{conceptual reframing} of fast-weight associative memories as world models
  (predictive substrates) rather than value functions, with a precise three-factor Hebbian
  instantiation.
\item A \textbf{seeded, multi-environment interception benchmark} with analytic, model-based,
  and model-free baselines, a reset-on-catch protocol that measures genuine interception
  rather than ``lock-on,'' and a forecast-error metric that isolates prediction quality.
\item \textbf{Empirical evidence} that (a) model-free value learning --- including a linear TD
  learner using the \emph{same} associative memory, a backprop DQN, and PPO --- fails
  catastrophically on interception, while (b) the same memory as a world model reaches
  near-optimal interception and (c) beats first- and second-order analytic extrapolation on
  curved prey.
\end{enumerate}

We are explicit about scope: this is a controlled toy domain, not a benchmark suite, and its
purpose is to isolate one conceptual question. Section~\ref{sec:discussion} discusses
limitations and what a fuller study would require.

\section{Related Work}\label{sec:related}

\textbf{Fast-weight memories.} The idea of using rapid weight changes as a form of short-term
memory traces back to \cite{schmidhuber1992}, and has been connected to linear attention
(``transformers are secretly fast-weight programmers'') \cite{schlag2021}. The Dragon Hatchling
\cite{bdh2025} centers a fast-weight associative memory as the bridge between
transformer-style attention and brain-like computation. Our work uses the core primitive ---
the outer-product weight update and similarity-weighted readout --- and asks what it should
represent.

\textbf{Hebbian and three-factor learning.} Classical Hebbian learning co-activates pre- and
post-synaptic units ($\Delta W \propto y x^\top$). Three-factor rules gate this plasticity by a
third, neuromodulatory signal (e.g., reward or prediction error), yielding biologically
plausible and provably convergent variants \cite{fremaux2016}. The delta (Widrow--Hoff / LMS)
rule $W \leftarrow W + \eta\,\delta\,x^\top$, where $\delta$ is the prediction error, is
precisely such a three-factor rule \cite{widrow1960}; we use its normalized form.

\textbf{Model-based RL and world models.} Learning a model of the environment and planning with
it is a classical RL paradigm. Dreamer \cite{hafner2023} popularized the split we instantiate:
train a world model densely on prediction error, then plan (imagine) rollouts with it, while
the policy is driven by sparse task reward. Our world model is the same idea at the scale of a
single associative matrix rather than a deep recurrent network.

\textbf{The deadly triad.} \cite{sutton2018} identify three properties whose combination can
cause divergence in TD learning: function approximation, bootstrapping, and off-policy updates.
Our linear-Q baseline is a textbook instance of all three, and its failure is a live
illustration.

\textbf{Pursuit and interception.} Interception is a classic control/geometry problem with
analytic solutions (lead pursuit, proportional navigation) \cite{shneydor1998,nahin2007} and
biological relevance \cite{fajen2003}. We use it as a clean, interpretable test bed rather than
contributing new pursuit algorithms.

\section{Problem Setting}\label{sec:problem}

\subsection{Environment}

We use a two-dimensional toroidal world of size $1200 \times 800$ px with discrete-time
dynamics at $dt = 0.05$ s. A single \textbf{prey} moves autonomously; a single \textbf{chaser}
(pursuer) attempts to catch it. A catch occurs when the chaser comes within a capture radius of
$48$ px; a $0.5$ s cooldown follows before the next catch can register.

\begin{itemize}
\item \textbf{Chaser.} Constant top speed $v_c = 175$ px/s. Its heading changes at a rate
  bounded by $\omega_{\max} = 3.6$ rad/s (a \emph{turn-rate clamp}). It cannot decelerate.
\item \textbf{Prey.} Speed $165$ px/s ($0.94\times$ the chaser), so the chaser cannot win by
  raw speed --- it must \emph{intercept}. Six dynamics (Table~\ref{tab:results}, top):
  \begin{itemize}
  \item \texttt{const-vel} --- constant velocity (linear; first-order extrapolation is exact).
  \item \texttt{circling} --- constant turn rate $\dot\theta \in \pm 1.2$ rad/s drawn per
    episode, so the prey runs in persistent arcs/circles.
  \item \texttt{ou-turn}, \texttt{ou-vel} --- Ornstein--Uhlenbeck (mean-reverting, stochastic)
    turn and speed.
  \item \texttt{jump} --- piecewise-constant heading with random jumps.
  \item \texttt{flee} --- reactive avoidance: steers away from the chaser and speeds up when
    approached.
  \end{itemize}
\end{itemize}

The first two are the scientifically decisive pair: \texttt{const-vel} is where first-order
extrapolation is already Bayes-optimal (a control), and \texttt{circling} is where intercepting
a \emph{trajectory} --- not a point velocity --- is what separates predictors. The stochastic
(\texttt{ou-*}, \texttt{jump}) and reactive (\texttt{flee}) prey test robustness and are where
we document honest limits.

\subsection{Protocol and metrics}

A naive ``catches per fixed horizon'' metric is dominated by \textbf{lock-on}: once a faster
chaser is within the capture radius it rarely loses the prey, so every competent agent
saturates at the cooldown rate and prediction quality is invisible. We therefore use a
\textbf{reset-on-catch} protocol: on each catch the prey is teleported to a uniformly random
location 400--600 px from the chaser. Every catch is then a \emph{fresh interception from
distance}, and the catch count directly measures interception efficiency.

We report \textbf{catches per episode} (mean $\pm$ standard deviation over seeds) as the
primary metric, and a \textbf{forecast error} --- mean absolute error between a predictor's
forecast of the prey's future position and its true position, evaluated at a fixed 1--2 s
horizon --- as a direct measure of prediction quality independent of the catch/steering loop.

\subsection{The oracle ceiling}

Interception has a clean ceiling. If the chaser knew the prey's future trajectory exactly and
could steer optimally under the turn-rate clamp, it would achieve an optimal catch rate. In
practice the relevant ceiling for our comparison is the \textbf{lead-pursuit ceiling} --- the
catch rate of a chaser that steers perfectly at the true future prey position --- which the
analytic baselines approximate from below. We do not claim the chaser reaches the true optimum;
we claim it reaches the lead-pursuit regime, and that the world model does so \emph{without} a
hand-crafted dynamics model.

\section{Method}\label{sec:method}

\subsection{The fast-weight associative memory}

Let $s \in \mathbb{R}^{d}$ be a normalized sensory state and $y \in \mathbb{R}^{2}$ a
two-dimensional target (a velocity, below). A fast-weight memory is a matrix
$W \in \mathbb{R}^{2 \times d}$ of associations, written and read as:

\begin{equation}
\text{read:}\quad \hat{y} = W \phi(s), \qquad
\text{write:}\quad W \leftarrow W + \eta\,\underbrace{(y - \hat{y})}_{\delta}\,\phi(s)^{\top}
\Big/ \big(1 + \|\phi(s)\|^2\big).
\end{equation}

The write is a \textbf{three-factor (error-gated) Hebbian} rule: the outer product of a
pre-synaptic activity $\phi(s)$ and a post-synaptic prediction error $\delta$, normalized for
stability (normalized LMS), with a small synaptic weight decay. The read is
\textbf{content-addressable}: $\hat y = W\phi(s)$ is the similarity-weighted recall of all
stored associations, so the memory \emph{generalizes across states} rather than storing a
lookup table. This is the BDH fast-weight primitive; nothing else is added.

The sensory features $\phi(s)$ are, by default, \textbf{linear}:
$\phi(s) = [\,1,\ \Delta x/R,\ \Delta y/R,\ v_x/v_{\max},\ v_y/v_{\max}\,]$, where $\Delta$ is
the wrapped relative position of the prey and $v$ its velocity. We deliberately do \emph{not}
hand the model the prey's turn rate or acceleration; it must infer curvature from raw
observations. (An optional random-Fourier expansion can be added to make the readout nonlinear;
Section~\ref{sec:ablations} shows it is unnecessary here.)

\subsection{The world model}

We use the memory as a \textbf{world model}: it learns the prey's one-step transition
$f: s_t \mapsto v_{t+1}$, mapping the current state to the \emph{next} velocity. Training is
dense and self-supervised --- every observed tick updates the memory toward the true next
velocity --- with no reward signal. To predict the prey's position $\tau$ steps ahead, the model
is \textbf{rolled forward (imagined)}: starting from the current state, it repeatedly predicts
the next velocity and integrates it to a new position, for $\tau$ steps.

\subsection{The planner}

Planning is deliberately minimal and shared across all predictors: \textbf{lead pursuit} ---
steer at the predicted future prey position. The lead horizon is $\tau = d / v_c$ (the time to
reach the prey at current speed), capped at $2$ s, so the chaser aims at where the prey
\emph{will be}, not where it is. The chaser's heading is then moved toward that aim under the
turn-rate clamp. The policy (which way to steer) is thus driven only by the forecast; the only
thing that varies across predictors is \emph{how the forecast is made}.

\subsection{The contrast: the same memory as a value function}

To make the value-vs-world-model contrast clean, we implement the evaluative reading with the
\emph{identical} substrate: $Q(s,a) = w_a \cdot \phi(s)$ for $8$ discrete headings, trained by
TD(0) \cite{sutton1988} with $\epsilon$-greedy exploration (Section~\ref{sec:experiments}).
This is exactly a linear value function over the same features $\phi$, updated by a
bootstrapped, off-policy rule --- the textbook ``deadly triad'' configuration
\cite{sutton2018}. Any performance gap between it and the world model is therefore attributable
to the \emph{role} of the memory (evaluate vs.\ predict), not to the substrate, features, or
representational capacity.

\section{Experiments}\label{sec:experiments}

\subsection{Setup}

All results are seeded and reproducible; each cell is a mean over seeds (10 for predictors, 3
for the slower model-free learners) with standard deviation. Hyperparameters were not tuned
per-environment except where an ablation explicitly varies them. The BDH world model uses the
linear feature set (Section~\ref{sec:method}), learning rate $\eta = 0.1$, and weight decay
$\lambda = 10^{-3}$.

\textbf{Baselines.}
\begin{itemize}
\item \emph{Predictors} (lead pursuit, differing only in the forecast):
  \begin{itemize}
  \item \texttt{velocity-lead}: $\hat p = p + v\tau$ (first-order extrapolation).
  \item \texttt{accel-lead}: $\hat p = p + v\tau + \tfrac12 a \tau^2$ (second-order Taylor,
    $a$ from a finite-difference estimate).
  \item \texttt{kalman-lead}: a constant-velocity Kalman filter + lead.
  \item \texttt{bdh}: the fast-weight world model (Section~\ref{sec:method}).
  \end{itemize}
\item \emph{Policies} (learn or select the heading directly):
  \begin{itemize}
  \item \texttt{pure-pursuit}: steer at the \emph{current} prey position (the no-prediction
    reflex).
  \item \texttt{mpc}: model-predictive control --- search $8$ headings, simulate a short
    horizon with a first-order (constant-velocity) prey model, minimize final distance.
  \item \texttt{linear-q}: linear TD Q-learning over the same features $\phi$ (the evaluative
    reading).
  \item \texttt{dqn}: an MLP Q-network with replay buffer and target network.
  \item \texttt{ppo}: clipped-surrogate PPO with policy and value MLPs.
  \end{itemize}
\end{itemize}

The model-free learners receive a dense reward $r = -\text{distance}/600$ plus a $+1$ catch
bonus, so their failure is a credit-assignment/interception failure rather than an artifact of
sparse reward.

\subsection{Result 1: world model vs.\ value function}

Table~\ref{tab:results} (bottom) reports catches under the reset-on-catch protocol. The
contrast is stark and consistent across all six environments:

\begin{itemize}
\item The \textbf{linear-Q learner --- the same associative memory used as a value function ---
  collapses} to $0$--$27$ catches, roughly $2$--$10\%$ of the lead-pursuit ceiling.
\item The \textbf{backprop DQN and PPO} do not exceed the \emph{no-prediction reflex}: on most
  environments they are \emph{worse} than pure pursuit, which does no learning and no
  prediction at all.
\item The \textbf{BDH world model}, by contrast, reaches $165$--$384$ catches ---
  $\sim\!94\%$ of the lead-pursuit ceiling --- with a fraction of the training and no value
  function.
\end{itemize}

Because linear-Q and the world model use the \emph{same} features and the \emph{same}
associative substrate, the gap is clean evidence that the fast-weight memory succeeds when it
\emph{predicts} and fails when it \emph{evaluates}.

\begin{table}[t]
\centering
\caption{Catches per episode (mean $\pm$ sd), reset-on-catch. \emph{Top:} lead-pursuit
predictors (10 seeds $\times$ 24{,}000 steps). \emph{Bottom:} policies (3 seeds $\times$
20{,}000 steps).}
\label{tab:results}
\small
\begin{tabular}{lcccc}
\toprule
prey & velocity-lead & accel-lead & kalman-lead & \textbf{BDH (world model)} \\
\midrule
const-vel & $364 \pm 57$ & $364 \pm 57$ & $148 \pm 50$ & $\mathbf{363 \pm 53}$ \\
circling  & $327 \pm 35$ & $131 \pm 192$ & $35 \pm 66$ & $\mathbf{384 \pm 38}$ \\
ou-turn   & $368 \pm 17$ & $364 \pm 10$ & $116 \pm 7$ & $330 \pm 11$ \\
ou-vel    & $388 \pm 8$ & $385 \pm 15$ & $136 \pm 16$ & $364 \pm 12$ \\
jump      & $342 \pm 10$ & $293 \pm 8$ & $133 \pm 10$ & $298 \pm 11$ \\
flee      & $182 \pm 2$ & $69 \pm 7$ & $6 \pm 3$ & $165 \pm 5$ \\
\bottomrule
\end{tabular}

\bigskip

\begin{tabular}{lccccc}
\toprule
prey & pure-pursuit & mpc & linear-q & dqn & ppo \\
\midrule
const-vel & $60 \pm 10$ & $242 \pm 33$ & $19 \pm 6$ & $30 \pm 28$ & $18 \pm 5$ \\
circling  & $179 \pm 52$ & $254 \pm 42$ & $22 \pm 5$ & $41 \pm 5$ & $23 \pm 5$ \\
ou-turn   & $129 \pm 10$ & $286 \pm 23$ & $25 \pm 6$ & $29 \pm 4$ & $23 \pm 1$ \\
ou-vel    & $137 \pm 11$ & $304 \pm 10$ & $27 \pm 5$ & $29 \pm 3$ & $21 \pm 9$ \\
jump      & $254 \pm 8$ & $265 \pm 11$ & $20 \pm 2$ & $44 \pm 9$ & $23 \pm 7$ \\
flee      & $160 \pm 1$ & $140 \pm 3$ & $0 \pm 1$ & $37 \pm 9$ & $8 \pm 8$ \\
\bottomrule
\end{tabular}
\end{table}

The clean ordering across approaches is \emph{pure-pursuit (reflex) < MPC (first-order search)
< velocity-lead (first-order, continuous) < BDH (learned world model)} on five of six
environments, with model-free value learning (linear-Q, DQN, PPO) \emph{below the reflex}
everywhere. Prediction and a learned model both matter, and planning with a model decisively
beats learning a value function. (On \texttt{flee}, where a first-order model is wrong, MPC
falls below the reflex --- the \emph{model}, not the \emph{planner}, is the limiting factor.)

\subsection{Result 2: world model vs.\ analytic extrapolation on curved prey}

On \texttt{circling} (Table~\ref{tab:results}, top), the BDH world model achieves
$384 \pm 38$ catches, \textbf{$+17\%$} over first-order extrapolation ($327 \pm 35$) and
\textbf{$+193\%$} over second-order ($131 \pm 192$). The second-order baseline is also wildly
unreliable --- its standard deviation ($192$) is larger than its mean --- whereas the world
model is stable ($\pm 38$).

The mechanism is visible in the forecast error: on \texttt{circling} the world model's
lead-horizon forecast error is $\sim\!62$--$69$ px versus $\sim\!81$--$100$ px for first-order.
Rolling forward the learned rotation tracks the arc; a straight-line extrapolation (first-order)
misses it, and a parabola (second-order) \emph{overshoots} it over the $\sim\!2$ s lead horizon
--- the $\tfrac12 a \tau^2$ term of a second-order Taylor expansion diverges from a circle. A
learned model that re-applies the true (learned) one-step dynamics at each of $\tau$ steps is
the correct object, and the fast-weight memory discovers it from raw observations without being
told the prey is circling.

On \texttt{const-vel} the world model ties first-order (as it should: the dynamics are linear,
and a linear model reproduces them exactly), confirming the model adds no spurious advantage. On
the stochastic prey (\texttt{ou-*}, \texttt{jump}) no predictor beats first-order --- the future
is genuinely unpredictable from the current state --- and the learned model adds a small amount
of variance. These are the honest negative results that delimit the claim.

\subsection{Ablations}\label{sec:ablations}

We verify the result is not fragile. (All ablations on \texttt{circling}, reset-on-catch.)

\begin{itemize}
\item \textbf{Turn-rate clamp} ($1.8 \to 5.4$ rad/s): BDH beats first-order at \emph{every}
  turn rate ($262/288/306/312$ vs.\ $219/273/280/283$), so the advantage is not an artifact of
  the clamp's value.
\item \textbf{Speed ratio} (prey $150/165/180$ vs.\ chaser $175$): BDH is robust
  ($322/306/294$). Second-order only competes at slow prey ($336$ at $150$) and collapses at
  $165+$ ($140$) as its parabola overshoots.
\item \textbf{Memory width (Fourier features $M$)} ($0/8/24/48$): the \emph{linear} associator
  ($M=0$) is optimal ($306$); the nonlinear expansion only adds variance ($253$--$304$). The
  relevant dynamics here are smooth, and a linear content-addressable memory captures them.
\item \textbf{Learning rate $\eta$} ($0.1/0.3/0.5/1.0$): slow learning is best
  ($\eta=0.1 \to 319$); \emph{every} $\eta$ beats first-order ($280$).
\item \textbf{Weight decay $\lambda$} ($0/10^{-4}/10^{-3}/10^{-2}$): moderate decay is best
  ($10^{-3} \to 306$); every $\lambda$ beats first-order.
\item \textbf{Action discretization} ($4/8/16$ headings, model-free): linear-Q ($\sim\!22$),
  DQN ($\sim\!37$--$46$), and PPO ($\sim\!23$--$30$) stay far below the reflex ($\sim\!179$)
  at every discretization, so the model-free failure is \emph{not} a discretization artifact
  --- it persists at 16 actions.
\end{itemize}

\section{Discussion}\label{sec:discussion}

\subsection{Why model-free value learning fails here}

Three compounding factors explain the model-free collapse. First, the \textbf{deadly triad}
(linear-Q is linear function approximation + TD bootstrapping + $\epsilon$-greedy off-policy
updates \cite{sutton2018}) makes the value estimate unstable or divergent. Second, the
\textbf{turn-rate clamp} means the greedy action cannot be executed instantaneously: even a
perfect value function's argmax heading takes $\sim\!0.1$ s to reach, so one-step-greedy value
learning systematically under-turns. Third, the \textbf{interception reward is inherently
temporal}: catching a moving target requires anticipating a future position, which is a
prediction problem, not a credit-assignment problem. A value function must propagate reward
through many time steps of a high-frequency control loop to \emph{implicitly} represent the same
thing a world model \emph{explicitly} predicts in one step.

The dense reward and strong function approximators (DQN, PPO) do not rescue this: they reliably
learn to \emph{approach} the prey but not to \emph{intercept} it, and frequently do worse than
the reflex.

\subsection{Why the world model works}

Interception is a prediction problem. The world model's inductive bias --- learn the map
$s_t \mapsto v_{t+1}$ and roll it forward --- is exactly aligned with the task, and its
training signal (dense, self-supervised prediction error) is abundant at every tick, independent
of the sparse catch reward. This is the Dreamer split \cite{hafner2023} in miniature: the policy
(steer at the forecast) is driven by outcome, while the model is driven by prediction. The fact
that a \emph{linear} associative memory suffices here suggests that, for smooth physical
dynamics, content-addressable recall of local transitions is already a strong world model ---
and that the ``fast weights'' of the BDH architecture are best read as a predictive substrate.

\subsection{Limitations}

This is a controlled toy domain and a first step, not a claim of generality. Specifically:
\begin{enumerate}
\item The result is strongest where it should be --- curved, predictable motion --- and does
  not hold on genuinely unpredictable (Martingale) prey, where no predictor beats first-order.
\item On \textbf{reactive} prey (\texttt{flee}), the world model is only comparable to
  first-order; our imagination holds the chaser fixed while rolling the prey forward, which is
  a crude approximation of the closed-loop pursuit--evasion dynamics.
\item The world model is \textbf{linear}; the nonlinear (Fourier) expansion did not help on
  these smooth dynamics, so we do not demonstrate nonlinear world modeling.
\item The environment is single-prey, toroidal, and two-dimensional; real pursuit adds walls,
  multiple prey, partial observation, and sensing noise.
\item Baselines SAC (continuous control) and a circle-fit analytic predictor remain for future
  work, as do comparisons on a standard RL benchmark rather than a purpose-built task.
\end{enumerate}

\subsection{What would make this a full paper}

To elevate this from a technical report to a full paper we would: (i) fix the closed-loop
imagination so reactive prey are handled; (ii) add SAC and a circle-fit (or Kalman-with-turn)
analytic baseline to close the baseline set; (iii) sweep stochasticity and noise to establish
the boundary of ``predictable enough''; and (iv) reproduce the core value-vs-world-model
contrast on a standard interception benchmark with multiple seeds and significance tests. The
conceptual claim --- fast weights predict, not evaluate --- does not depend on these, but its
generality does.

\section{Conclusion}\label{sec:conclusion}

We asked whether the fast-weight associative memory at the heart of the Dragon Hatchling
architecture is a value function or a world model. In a seeded interception benchmark, the
answer is decisive: the \emph{same} memory, used to predict, reaches $\sim\!94\%$ of the
lead-pursuit ceiling and beats first- and second-order analytic extrapolation on curved prey,
while used to evaluate it collapses below a no-prediction reflex --- and even strong model-free
learners (DQN, PPO) cannot match the reflex. The result supports a clean reading of fast-weight
memories as predictive substrates trained by dense self-supervised error, with sparse
outcome-driven planning on top --- a Dreamer-style division of labor instantiated by a
biologically plausible three-factor Hebbian rule. We hope this reframing is useful both to
readers of the BDH literature and to practitioners choosing what to put in a fast-weight
memory: put the \emph{next state}, not the \emph{value}.

\section*{Acknowledgments and disclosures}

\textbf{Author responsibilities.} All authors take full responsibility for the content of this
manuscript, regardless of how it was produced.

\textbf{Use of generative AI language tools.} Portions of this manuscript were drafted and
edited with the assistance of a text-to-text generative AI language model. All results, code,
references, and claims were reviewed and verified by the authors, who accept full
responsibility for them; the AI tool is not an author. This statement is included to comply
with arXiv's policy on authors' use of generative AI language tools
(\url{https://info.arxiv.org/help/moderation/index.html}).

\textbf{Code and data availability.} The benchmark, all experiment configurations, and the code
to reproduce every number in this paper are available at
\url{https://github.com/llvm-x86/fast-weights-predict}.

\begin{thebibliography}{11}

\bibitem{bdh2025}
A.~Kosowski, P.~Uzna\'nski, J.~Chorowski, Z.~Stamirowska, and M.~Bartoszkiewicz.
\newblock The Dragon Hatchling: The Missing Link between the Transformer and Models of the
  Brain.
\newblock arXiv:2509.26507, 2025.

\bibitem{schmidhuber1992}
J.~Schmidhuber.
\newblock Learning to control fast-weight memories: An alternative to dynamic recurrent
  networks.
\newblock \emph{Neural Computation}, 4(1):131--139, 1992.

\bibitem{schlag2021}
I.~Schlag, K.~Irie, and J.~Schmidhuber.
\newblock Linear transformers are secretly fast weight programmers.
\newblock In \emph{Proceedings of the 38th International Conference on Machine Learning
  (ICML)}, 2021.

\bibitem{sutton2018}
R.~S. Sutton and A.~G. Barto.
\newblock \emph{Reinforcement Learning: An Introduction}.
\newblock MIT Press, 2nd edition, 2018.

\bibitem{sutton1988}
R.~S. Sutton.
\newblock Learning to predict by the methods of temporal differences.
\newblock \emph{Machine Learning}, 3(1):9--44, 1988.

\bibitem{fremaux2016}
N.~Fr\'emaux and W.~Gerstner.
\newblock Neuromodulated spike-timing-dependent plasticity, and theory of three-factor learning
  rules.
\newblock \emph{Frontiers in Neural Circuits}, 9:85, 2016.

\bibitem{hafner2023}
D.~Hafner, J.~Pasukonis, J.~Ba, and T.~Lillicrap.
\newblock Mastering diverse domains through world models.
\newblock arXiv:2301.04104, 2023.

\bibitem{widrow1960}
B.~Widrow and M.~E. Hoff.
\newblock Adaptive switching circuits.
\newblock In \emph{IRE WESCON Convention Record}, 1960.

\bibitem{shneydor1998}
N.~A. Shneydor.
\newblock \emph{Missile Guidance and Pursuit: Kinematics, Dynamics and Control}.
\newblock Horwood Publishing, 1998.

\bibitem{nahin2007}
P.~J. Nahin.
\newblock \emph{Chases and Escapes: The Mathematics of Pursuit and Evasion}.
\newblock Princeton University Press, 2007.

\bibitem{fajen2003}
B.~R. Fajen and W.~H. Warren.
\newblock Behavioral dynamics of steering, obstacle avoidance, and route selection.
\newblock \emph{Journal of Experimental Psychology: Human Perception and Performance},
  29(2):343--362, 2003.

\end{thebibliography}

\end{document}
